\documentclass[11pt]{article}
\usepackage[utf8]{inputenc}
\usepackage[T1]{fontenc}
\usepackage{lmodern}
\usepackage[margin=0.9in]{geometry}
\usepackage{enumitem}
\usepackage{xcolor}
\usepackage{hyperref}
\usepackage{amsmath, amssymb}

\hypersetup{
  colorlinks = true,
  linkcolor = blue!50!black,
  urlcolor = blue!50!black
}

\setlist[itemize]{leftmargin=1.4em, itemsep=0.25em, topsep=0.35em}
\setlength{\parindent}{0pt}
\setlength{\parskip}{0.55em}
\newcommand{\scriptlabel}{\textbf{Script. }}
\newcommand{\sectionlabel}[1]{\textbf{#1}}
\newcommand{\actioncue}[1]{\textit{\textbf{#1}}}
\title{Supervised Bayesian Matrix Factorization\\[0.8em]\large Presentation Notes}
\author{Andrew Walther}
\date{April 9, 2026}

\begin{document}
\maketitle

\noindent This companion is written as a preparation guide for an approximately 15-minute lab meeting talk. The main deck sections are written as the primary script. The backup sections are shorter and are intended for discussion only if those slides come up.

\section*{Opening}

\subsection*{Title Slide}
\scriptlabel Thanks everyone for letting me share an update on the supervised Bayesian matrix factorization project. Today I want to show both the statistical idea and the first cross-cohort pancreatic ductal adenocarcinoma (PDAC) results. The short version is that bringing survival directly into the factorization changes which latent programs the model prioritizes, and that seems to matter in real data.

\begin{itemize}
\item \sectionlabel{Main Takeaway:} Frame the talk as both a methods update and a first results update.
\item \sectionlabel{Emphasis:} Set the audience expectation that the talk is about clinically relevant latent structure, not only reconstruction quality.
\item \sectionlabel{Timing:} Keep this to one or two sentences and move quickly into the outline.
\end{itemize}

\section*{Main Deck}

\subsection*{Slide 1: Outline}
\scriptlabel I will spend just a couple of minutes on the motivation, then walk through the model and inference at a high level, and spend most of the talk on the cross-dataset results. I will end with the key takeaways and what I think the next practical steps are for turning the latent programs into something biologically interpretable.

\begin{itemize}
\item \sectionlabel{Main Takeaway:} This is a results-forward talk, not a derivation-heavy talk.
\item \sectionlabel{Emphasis:} Mention that most of the time will be spent on the prior-by-$K$ comparison and the cohort deep dives.
\end{itemize}

\subsection*{Slide 2: The PDAC Problem}
\scriptlabel Pancreatic ductal adenocarcinoma, or PDAC, remains one of the deadliest solid tumours, and unlike some other cancers we still do not have widely adopted molecular stratification at diagnosis. The problem is not that gene expression lacks signal, but that the signal is mixed across malignant, stromal, and immune sources. The goal here is to identify gene expression programs whose patient-level activation is not only biologically present, but also prognostically informative.

\begin{itemize}
\item \sectionlabel{Main Takeaway:} There is a real clinical need for prognostic molecular structure in PDAC.
\item \sectionlabel{Figure/Table Purpose:} The callout translates the statistical objective into clinical use cases such as trial stratification and targeted therapy.
\item \sectionlabel{Emphasis:} Distinguish ``describing tumour heterogeneity'' from ``predicting disease course.''
\end{itemize}

\subsection*{Slide 3: Why Supervised $>$ Unsupervised}
\textbf{Script.} The standard workflow is to run PCA or NMF first, then test the derived factors for association with survival. The weakness is that unsupervised methods rank factors by expression variance, not by clinical relevance. A factor can explain a lot of variance and still be biologically uninteresting for prognosis, while a smaller-variance factor can be the one that actually separates patients by outcome. Supervision changes the objective so the latent space has to serve both expression reconstruction and survival prediction.

\begin{itemize}
\item \sectionlabel{Main Takeaway:} Post hoc testing after PCA can miss clinically useful structure.
\item \sectionlabel{Emphasis:} The key conceptual win is not just prediction, but separation of prognostic from non-prognostic latent structure.
\item \sectionlabel{Good Line To Stress:} If $\hat{\beta}_k = 0$, the model keeps the factor for expression but not for prognosis.
\end{itemize}

\subsection*{Slide 4: The Gap \& Our Approach}
\textbf{Script.} Existing supervised matrix factorization approaches usually rely on fixed penalties and manual tuning to balance reconstruction against survival. What we want instead is a framework where shrinkage adapts to the data, uncertainty is carried through the updates, and the model can tell us when five factors are simply not enough. That is the niche SBMF is trying to fill.

\begin{itemize}
\item \sectionlabel{Main Takeaway:} SBMF is meant to improve on fixed-penalty supervised factorization, not replace the underlying idea of supervision.
\item \sectionlabel{Figure/Table Purpose:} The callout anchors the comparison against penalised approaches such as deSurv without making the slide about that comparison alone.
\item \sectionlabel{Emphasis:} The ELBO gives a natural balance, so there is no external $\alpha$ to tune between genomics and survival.
\end{itemize}

\subsection*{Slide 5: The SBMF Model}
\textbf{Script.} This slide is the core model statement. The expression matrix is written as a low-rank factorization with gene-specific noise precision, and the exact same latent patient loadings feed the Cox proportional hazards model. That shared latent space is the whole point: each factor is encouraged to explain expression structure while also being allowed to matter, or not matter, for survival through its corresponding $\beta_k$.

\begin{itemize}
\item \sectionlabel{Main Takeaway:} One latent space drives both the genomics and survival components.
\item \sectionlabel{Figure/Table Purpose:} The table defines the objects once so the later slides can move faster.
\item \sectionlabel{Emphasis:} Each column of $F$ is a gene program, each row of $L$ is a patient’s program activation profile.
\end{itemize}

\subsection*{Slide 6: Dual-Source Objective Function}
\textbf{Script.} The optimisation target is the ELBO, which can be read as expression fit plus survival fit minus regularisation. The important interpretation is that the shared $L$ appears in both likelihood terms, so the latent loadings cannot specialise to only one task. This is how the supervision is built into the factorization itself rather than added afterward.

\begin{itemize}
\item \sectionlabel{Main Takeaway:} Supervision is implemented through the joint variational objective.
\item \sectionlabel{Figure/Table Purpose:} The left and right callouts separate the two likelihood pieces so the audience can see how the ELBO is assembled.
\item \sectionlabel{Emphasis:} Say ``both likelihoods share the same latent loadings'' explicitly.
\end{itemize}

\subsection*{Slide 7: EBNM \& Adaptive Shrinkage}
\textbf{Script.} The easiest way to explain EBNM is to start from ridge and LASSO. In penalised regression, we choose a fixed penalty strength. Here, the prior itself is estimated from the data, so shrinkage can adapt to what the current factor or coefficient estimates look like. I tested both point-normal and point-laplace families because they give different tail behaviour and sparsity patterns.

\begin{itemize}
\item \sectionlabel{Main Takeaway:} The prior is learned from the data rather than fixed in advance.
\item \sectionlabel{Figure/Table Purpose:} The right column gives the practical interpretation of the two prior families that will later appear in the results table.
\item \sectionlabel{Emphasis:} Mention the spike-and-slab intuition and the fact that posterior second moments are carried forward.
\end{itemize}

\subsection*{Slide 8: Inference: Factor-Wise CAVI}
\textbf{Script.} Inference is done with Gauss-Seidel coordinate ascent variational inference. For each factor, we update the patient loadings, the gene weights, and then the survival coefficient, each through an EBNM subproblem. The automatic $K$ step is currently a pragmatic screen rather than full cross-validation: after fitting a large model, we keep factors that either explain non-trivial variance or carry non-negligible survival signal.

\begin{itemize}
\item \sectionlabel{Main Takeaway:} The algorithm updates one factor at a time and then screens active factors from a larger fit.
\item \sectionlabel{Figure/Table Purpose:} The right column explains both convergence and the current $K_{\text{eff}}$ rule.
\item \sectionlabel{Emphasis:} Say clearly that CV-based $K$ selection is still a next step; today’s $K_{\text{eff}}$ is a useful heuristic, not the final model-selection solution.
\end{itemize}

\subsection*{Slide 9: Data Intro --- PDAC Cohorts}
\textbf{Script.} Before showing results, I want to make the validation setting concrete. These are seven independent PDAC cohorts spanning RNA-seq, microarray, and proteomics, with sample sizes from 52 to 288. That breadth matters because a method can look good on one cohort and one assay, but real robustness means behaving reasonably across platforms with different noise structures.

\begin{itemize}
\item \sectionlabel{Main Takeaway:} The validation is cross-cohort and cross-platform, not a single-dataset demonstration.
\item \sectionlabel{Figure/Table Purpose:} The table is there to ground the audience in platform, cohort size, event counts, and censoring at a glance.
\item \sectionlabel{Emphasis:} All analyses use the top 5,000 most-variable features and column centring for comparability.
\end{itemize}

\subsection*{Slide 10: Synthetic Validation}
\textbf{Script.} I start with synthetic data because there the ground truth is known. The model recovers the latent factors up to permutation, recovers the survival coefficients with the right sign and magnitude, and correctly drives the null factor’s coefficient to zero. That is the minimum sanity check before trusting any real-data story.

\begin{itemize}
\item \sectionlabel{Main Takeaway:} The model can recover known factor structure and a true null survival effect.
\item \sectionlabel{Figure/Table Purpose:} The left panel shows factor correspondence; the right panel shows survival-coefficient recovery.
\item \sectionlabel{Emphasis:} The null factor is the internal negative control for whether the shrinkage prior is doing something sensible.
\end{itemize}

\subsection*{Slide 11: Convergence \& Baseline Performance}
\textbf{Script.} Across all seven cohorts, the optimisation behaves stably. The RMSE traces descend cleanly, although the number of iterations varies a lot by platform, with proteomics taking longer than the microarray cohorts. The baseline point-normal, $K=5$ configuration is not the best model; I show it here mainly to establish that the optimiser converges and to motivate why the later prior-by-$K$ comparison matters.

\begin{itemize}
\item \sectionlabel{Main Takeaway:} Convergence is reliable across cohorts, but the default baseline is clearly not sufficient.
\item \sectionlabel{Figure/Table Purpose:} The two traces contrast a large microarray cohort against a noisier proteomics cohort.
\item \sectionlabel{Emphasis:} Note that platform differences show up in iteration count, which is exactly what one would expect if noise structures differ.
\end{itemize}

\subsection*{Slide 12: Prior $\times$ $K$ Factorial: The Central Result}
\textbf{Script.} This is the main results table. Each cohort is evaluated under four conditions: point-normal and point-laplace, each at fixed $K=5$ and at data-adapted $K_{\text{eff}}$. The two patterns I want the audience to see are that point-normal at $K=5$ often degenerates badly, and that allowing the model to use more factors improves performance substantially overall, especially for the point-normal fits.

\begin{itemize}
\item \sectionlabel{Main Takeaway:} Fixed $K=5$ is too restrictive, and the point-normal baseline fails in several cohorts.
\item \sectionlabel{Figure/Table Purpose:} The table shows the full factorial comparison without hiding any cohort-specific variation.
\item \sectionlabel{Emphasis:} Be precise. Say ``mean performance improves markedly under $K_{\text{eff}}$'' rather than ``every condition improves in every cohort.''
\end{itemize}

\subsection*{Slide 13: Prior $\times$ $K$ Scatter Plot}
\textbf{Script.} The scatter plot is the visual version of the previous table. The blue points show the unstable point-normal, $K=5$ fits. Moving to point-laplace helps, but the strongest overall shift is what happens when the model is allowed to move to $K_{\text{eff}}$, which pulls most cohorts upward.

\begin{itemize}
\item \sectionlabel{Main Takeaway:} The visual trend matches the table, with the biggest rescue coming from data-adapted $K$.
\item \sectionlabel{Figure/Table Purpose:} This slide compresses the same factorial result into an easier pattern-recognition view.
\item \sectionlabel{Emphasis:} Use this slide for trend language, not for over-detailed numeric discussion.
\end{itemize}

\subsection*{Slide 14: Key Findings: Prior $\times$ $K$ Comparison}
\textbf{Script.} The summary result is that point-laplace is clearly better when $K$ is artificially fixed at 5, but once $K$ is allowed to adapt, the prior family matters much less than the number of factors. In other words, I would rather spend effort getting $K$ right than spending that same effort debating point-normal versus point-laplace. The scree plot on the right reinforces that the data often wants more than five active factors.

\begin{itemize}
\item \sectionlabel{Main Takeaway:} $K$ selection is the dominant lever in this analysis.
\item \sectionlabel{Figure/Table Purpose:} The scree plot gives a concrete example of why $K=5$ was too conservative.
\item \sectionlabel{Emphasis:} Three cohorts hit $K_{\max}=10$, which suggests the latent structure may be richer than the current cap allows.
\end{itemize}

\subsection*{Slide 15: Deep Dive --- Puleo Array}
\textbf{Script.} Puleo is the clearest real-data example because it is large and shows well-separated factor behaviour. The key observation is that factor 2 explains a lot of expression variance but is shrunk to $\hat{\beta}=0$, while other factors carry strong prognostic signal. That is exactly the separation we want from supervised factorization: not every large expression axis deserves to be treated as clinically important.

\begin{itemize}
\item \sectionlabel{Main Takeaway:} The model separates expression-dominant and prognosis-dominant programs in a large cohort.
\item \sectionlabel{Figure/Table Purpose:} The left panel shows which factors are prognostic; the right panel shows that those factors translate into visible survival stratification.
\item \sectionlabel{Emphasis:} Explicitly contrast this with what a post hoc PCA-based workflow would do.
\end{itemize}

\subsection*{Slide 16: Deep Dive --- TCGA\_PAAD}
\textbf{Script.} I included TCGA because it is the benchmark dataset most people in the room will recognise. The same qualitative pattern shows up here: some factors clearly matter for prognosis and others do not, and moving from $K=5$ to $K_{\text{eff}}=9$ gives one of the biggest performance gains in the study. So the main story is not limited to Puleo.

\begin{itemize}
\item \sectionlabel{Main Takeaway:} The TCGA cohort reproduces the same supervised-factorization pattern in a familiar benchmark dataset.
\item \sectionlabel{Figure/Table Purpose:} The panels mirror the Puleo slide so the audience can compare cohort-specific structure directly.
\item \sectionlabel{Emphasis:} Say that the large improvement here is one reason to take $K$ mis-specification seriously.
\end{itemize}

\subsection*{Slide 17: Hold-Out Prediction}
\scriptlabel This slide asks the practical question: does the model help on unseen patients? In five of the seven cohorts, the best SBMF condition beats the PCA baseline on the hold-out set. I also want to be honest about the caveat that the smallest cohorts have very small test sets, so the C-index there is noisier and should not be over-interpreted.

\begin{itemize}
\item \sectionlabel{Main Takeaway:} SBMF usually generalises better than the PCA baseline, but the smallest cohorts remain unstable.
\item \sectionlabel{Figure/Table Purpose:} The table collapses the factorial comparison down to the best-performing condition per cohort.
\item \sectionlabel{Emphasis:} This is the slide to use for the ``does it help in prediction?'' question.
\end{itemize}

\subsection*{Slide 18: Future Directions \& Summary}
\scriptlabel The immediate next step is biological interpretation: the factors need names, pathways, and a clearer connection to known PDAC biology. On the methods side, the highest-value improvement is proper cross-validated $K$ selection on the cluster at a higher $K_{\max}$. The summary I want people to leave with is that the method converges, the supervision changes which latent programs matter, and getting $K$ right matters more than fine-tuning the prior family.

\begin{itemize}
\item \sectionlabel{Main Takeaway:} The method is working well enough to move toward biological interpretation and more principled model selection.
\item \sectionlabel{Emphasis:} Keep the summary crisp and repeat the ``$K$ matters most'' message one last time.
\item \sectionlabel{Good Closing Phrase:} The next stage is turning statistical factors into interpretable biology.
\end{itemize}

\subsection*{Slide 19: Discussion}
\scriptlabel \actioncue{Stop here, thank the group, and invite questions.} If discussion turns technical, I have backup slides on pooled analysis, EBNM details, convergence, heteroscedastic precision, and per-cohort diagnostics.

\begin{itemize}
\item \sectionlabel{Main Takeaway:} Open the room for discussion without adding new content.
\item \sectionlabel{Emphasis:} Mention the backup material so the audience knows deeper technical detail is available.
\end{itemize}

\section*{Backup Slides}

\subsection*{Backup 1: Pooled RNA-seq Analysis}
\scriptlabel If asked about shared signal across studies, I would use this slide to show that a pooled RNA-seq analysis is feasible once we restrict to genes shared across cohorts. The result is preliminary, but it suggests the model can be pushed toward cross-cohort latent programs rather than only within-cohort fits.

\begin{itemize}
\item \sectionlabel{Main Takeaway:} Pooled fitting is possible and useful for replication questions.
\item \sectionlabel{Figure Purpose:} The two panels show that the pooled fit both converges and yields interpretable survival coefficients.
\end{itemize}

\subsection*{Backup 2: EBNM Deep Dive}
\scriptlabel This is the technical explanation of the EBNM subproblem. I would only use it if someone asks what the updates are doing mathematically. The key point is that each update reduces to a noisy normal-means problem, and the posterior variance is explicitly retained through the second moment.

\begin{itemize}
\item \sectionlabel{Main Takeaway:} EBNM gives adaptive shrinkage and an uncertainty-aware second moment.
\item \sectionlabel{Emphasis:} The error-in-variables correction is the line to stress if discussion gets technical.
\end{itemize}

\subsection*{Backup 3: GEP Heatmap --- Puleo Array}
\scriptlabel This is a useful visual if someone wants to see what the latent programs look like across patients. I would use it to show that the programs induce visible patient structure rather than only producing scalar performance metrics.

\begin{itemize}
\item \sectionlabel{Main Takeaway:} The latent programs organize patients in structured, non-random ways.
\item \sectionlabel{Figure Purpose:} Show cohort heterogeneity in the loading matrix directly.
\end{itemize}

\subsection*{Backup 4: ELBO Convergence Traces}
\scriptlabel If someone asks whether the optimiser is well behaved beyond RMSE, this is the answer. The ELBO traces show the variational objective moving in the expected direction, which complements the reconstruction curves from the main talk.

\begin{itemize}
\item \sectionlabel{Main Takeaway:} Convergence is supported by both RMSE and variational-objective diagnostics.
\item \sectionlabel{Figure Purpose:} Provide a stronger optimisation sanity check than a single metric alone.
\end{itemize}

\subsection*{Backup 5: Tau/Precision Distribution}
\scriptlabel This slide is there for the heteroscedasticity question. The gene-specific precision estimates are clearly not constant, which justifies including $\tau_j$ rather than assuming one shared noise level for all features.

\begin{itemize}
\item \sectionlabel{Main Takeaway:} Real cohorts show feature-specific noise, so modelling $\tau_j$ is worthwhile.
\item \sectionlabel{Figure Purpose:} Highlight that the spread differs across platforms and cohorts.
\end{itemize}

\subsection*{Backup 6: Per-Cohort Beta Comparisons}
\scriptlabel If someone wants reassurance that the prognostic-factor pattern is not limited to Puleo and TCGA, this is the slide I would use. The exact factors differ by cohort, but the recurring theme is that some factors are clearly shrunk toward zero while others are retained as prognostic.

\begin{itemize}
\item \sectionlabel{Main Takeaway:} Prognostic versus non-prognostic separation appears across the cohort panel, not only in the highlighted examples.
\item \sectionlabel{Figure Purpose:} Cohort-by-cohort beta patterns for quick comparison.
\end{itemize}

\subsection*{Backup 7: Per-Cohort Kaplan-Meier Curves}
\scriptlabel This is the visual counterpart to the previous slide. I would use it if the discussion shifts toward clinical interpretability and patient stratification rather than model mechanics.

\begin{itemize}
\item \sectionlabel{Main Takeaway:} Many cohorts show meaningful survival separation when patients are grouped by program activation.
\item \sectionlabel{Figure Purpose:} Direct clinical-style visualisation of the model output.
\end{itemize}

\subsection*{Backup 8: Signal Recovery --- Synthetic Detail}
\scriptlabel This backup gives a closer look at the synthetic recovery result. It is useful if someone wants stronger evidence that the factor recovery is real and not only visible in a summary correlation matrix.

\begin{itemize}
\item \sectionlabel{Main Takeaway:} The simulation supports genuine recovery of latent structure.
\item \sectionlabel{Figure Purpose:} Point-by-point agreement between estimated and true loadings after sign correction.
\end{itemize}

\subsection*{Backup 9: Method Comparison}
\scriptlabel I would use this slide if the discussion becomes comparative and someone asks where SBMF sits relative to penalised supervised MF methods. The honest answer is that the comparison is conceptual at this stage: SBMF brings adaptive shrinkage, uncertainty-aware updates, and a direct Bayesian treatment of the joint objective.

\begin{itemize}
\item \sectionlabel{Main Takeaway:} SBMF differs mainly in adaptive priors, uncertainty propagation, and the way the joint objective is handled.
\item \sectionlabel{Figure Purpose:} Compact side-by-side positioning of the two modelling philosophies.
\item \sectionlabel{Emphasis:} Avoid overselling; the goal is to explain why the Bayesian route is attractive for this problem.
\end{itemize}

\end{document}
